\documentclass[UTF8]{article}
\usepackage{xeCJK}
\usepackage{amsmath}
\usepackage{booktabs}
\usepackage{graphicx}
\usepackage{hyperref}

\title{FinCast: 基于基础模型的金融时序预测文献综述}
\author{罗新锐}
\date{\today}

\begin{document}
\maketitle

\section{三大挑战}
金融时序预测在经济稳定、政策制定与投资决策中具有关键作用。然而，该领域长期面临三大核心挑战：

\begin{enumerate}
  \item \textbf{时间非平稳性（Temporal Non-stationarity）}：金融市场分布会随时间变化，受宏观经济、政策、投资者行为等因素影响。
  \item \textbf{多领域异质性（Multi-domain Diversity）}：股票、外汇、期货、加密货币等金融领域各自具有不同的模式与规律。
  \item \textbf{多时间分辨率（Multi-temporal Resolution）}：金融数据的时间尺度跨度大，从秒级高频交易到月度经济指标。
\end{enumerate}

传统模型（如ARIMA、GARCH）在非线性建模和模式迁移方面能力有限，而深度学习方法虽具强大表示力，但容易过拟合且泛化性弱。

\section{论文核心创新与贡献}
FinCast（Zhu et al., 2025）提出了首个针对金融时序的基础模型（foundation model），其创新点包括：

\begin{itemize}
  \item 引入 \textbf{Point-Quantile Loss (PQ-loss)}，联合优化点预测与分位数预测，提高模型在非平稳条件下的鲁棒性；
  \item 设计 \textbf{Token级稀疏专家混合（Sparse Mixture-of-Experts, MoE）} 机制，使模型能在不同金融领域中动态选择专家；
  \item 引入 \textbf{可学习频率嵌入（Learnable Frequency Embedding）}，用于编码时间分辨率信息，提升对多频率数据的适应能力；
  \item 在零样本（Zero-shot）预测中显著超越现有SOTA模型，平均MSE下降约20\%，MAE下降约23\%。
\end{itemize}

\section{FinCast 模型架构与关键组件}
FinCast 为 \textbf{decoder-only Transformer} 架构，主要由以下三个模块组成：

\subsection{输入模块（Input Tokenization Block）}
原始时间序列首先经过实例归一化（Instance Normalization）与残差多层感知机（Residual MLP）编码，随后加入频率嵌入：
\[
h_{input} = \text{ResidualMLP}(\tilde{X}) + \text{Embed}_{freq}(f)
\]
其中 $\text{Embed}_{freq}$ 为可学习的频率向量。

\subsection{解码器主体（Decoder MOE Backbone）}
采用因果自注意力（Causal Self-Attention）保证自回归一致性，每个token仅关注当前及过去信息：
\[
\text{Scores} = \text{softmax}(QK^\top + M)
\]
并通过稀疏专家层（Sparse MoE）实现动态路由与专家特化：
\[
\text{MoE}(h_n) = \sum_i g_{i,n} \cdot \text{MLP}_i(h_n)
\]

\subsection{输出模块（Output Block）}
输出层通过残差MLP映射至预测空间，并进行逆归一化以恢复原始尺度：
\[
Y = \hat{y} \cdot \sigma + \mu
\]

\section{与传统时序模型的对比}
表~\ref{tab:comparison} 对比了 FinCast 与常见时序预测模型（如 ARIMA、LSTM）在关键特性上的差异。

\begin{table}[h]
\centering
\caption{FinCast 与传统模型比较}
\label{tab:comparison}
\begin{tabular}{@{}llll@{}}
\toprule
模型 & 核心思想 & 优点 & 局限性 \\ \midrule
ARIMA & 线性自回归 & 解释性强 & 难处理非线性、非平稳数据 \\
LSTM & 序列记忆 & 捕获短期依赖 & 训练慢、长依赖衰减 \\
FinCast & 基础模型+MoE+PQ-loss & 泛化强、跨领域、抗噪声 & 训练成本高、需大规模数据 \\ \bottomrule
\end{tabular}
\end{table}

\section{个人思考与改进方向}
抱歉，学长，这里缺失可能是知识面还太窄了，并没有什么特别的idea
\section{结论}
FinCast 通过结合 Transformer、MoE 与 PQ-loss 的创新设计，实现了金融时序预测领域的重要突破。它不仅能在无微调的条件下跨领域泛化，还为构建通用金融AI模型提供了范式参考。

\begin{thebibliography}{9}
\bibitem{fincast}
Zhuohang Zhu, Haodong Chen, Qiang Qu, and Vera Chung.
\textit{FinCast: A Foundation Model for Financial Time-Series Forecasting.}
In Proceedings of the 34th ACM International Conference on Information and Knowledge Management (CIKM ’25), 2025.
\url{https://arxiv.org/abs/2508.19609v1}

\end{thebibliography}

\end{document}
